
\documentclass[main.tex]{subfiles}

\begin{document}
\graphicspath{{imagenes/tema2/}{../imagenes/tema2/}}

% Diapositiva original 1
\portadatema{2}{Cálculo diferencial. Aplicaciones (I)}
% Diapositiva original 2
\begin{frame}{Noción de derivada}
\textbf{Definición de derivada}\par
La derivada de $f$ en $x$ viene dada por la expresión
\[f'(x)=\lim_{\Delta x\to0}\frac{f(x+\Delta x)-f(x)}{\Delta x}.\]
Para todos los valores de $x$ para los que exista este límite, $f'$ es una función de $x$.
\medskip

\textbf{Notaciones}
\[f'(x),\quad\frac{dy}{dx},\quad y',\quad\frac{d}{dx}[f(x)],\quad D_x[y].\]
\end{frame}
% Diapositiva original 3
\begin{frame}{Reglas básicas de derivación}
\scriptsize\renewcommand{\arraystretch}{1.4}
\resizebox{\textwidth}{!}{\begin{tabular}{|l|l|l|l|}\hline
\multicolumn{4}{|c|}{Reglas generales de derivación}\\\hline
Producto por un número & $(cf)'=cf'$ & & \\
Suma & $(f+g)'=f'+g'$ & Diferencia & $(f-g)'=f'-g'$\\
Producto & $(fg)'=f'g+fg'$ & Cociente & $\left(\frac fg\right)'=\frac{f'g-fg'}{g^2}$\\\hline
\multicolumn{4}{|c|}{Derivadas de funciones algebraicas}\\\hline
Constante & $c'=0$ & Potencia & $(x^n)'=nx^{n-1}$, $x'=1$\\\hline
\multicolumn{4}{|c|}{Derivadas de funciones trigonométricas}\\\hline
Seno & $(\sin x)'=\cos x$ & Coseno & $(\cos x)'=-\sin x$\\
Tangente & $(\tan x)'=\sec^2x$ & Cotangente & $(\cot x)'=-\csc^2x$\\
Secante & $(\sec x)'=\sec x\tan x$ & Cosecante & $(\csc x)'=-\csc x\cot x$\\\hline
\multicolumn{4}{|c|}{Regla de la cadena}\\\hline
Cadena & $[f(u)]'=f'(u)u'$ & Potencia & $(u^n)'=nu^{n-1}u'$\\\hline
\end{tabular}}
\end{frame}
% Diapositiva original 4
\begin{frame}{Derivación implícita}
\textbf{Ejemplo}\par Obtener la derivada de $y$ respecto a $x$ de $y^3+y^2-5y-x^2=-4$.\medskip

\textbf{Paso 1)} Derivar ambos lados de la ecuación respecto de $x$.
\[\frac{d}{dx}[y^3+y^2-5y-x^2]=\frac{d}{dx}[-4].\]
\[\frac{d}{dx}[y^3]+\frac{d}{dx}[y^2]-\frac{d}{dx}[5y]-\frac{d}{dx}[x^2]=\frac{d}{dx}[-4].\]
\[3y^2\frac{dy}{dx}+2y\frac{dy}{dx}-5\frac{dy}{dx}-2x=0.\]
\end{frame}
% Diapositiva original 5
\begin{frame}{Derivación implícita}
\textbf{Ejemplo}\par Obtener la derivada de $y$ respecto a $x$ de $y^3+y^2-5y-x^2=-4$.\medskip

\textbf{Paso 2)} Agrupar todos los términos en los que aparezca $dy/dx$ en el lado izquierdo de la ecuación y pasar todos los demás a la derecha.
\[3y^2\frac{dy}{dx}+2y\frac{dy}{dx}-5\frac{dy}{dx}=2x.\]
\end{frame}
% Diapositiva original 6
\begin{frame}{Derivación implícita}
\textbf{Ejemplo}\par Obtener la derivada de $y$ respecto a $x$ de $y^3+y^2-5y-x^2=-4$.\medskip

\textbf{Paso 3)} Sacar factor común $dy/dx$ en el lado izquierdo de la ecuación.
\[\frac{dy}{dx}(3y^2+2y-5)=2x.\]
\textbf{Paso 4)} Despejar la derivada de $y$ respecto a $x$.
\[\frac{dy}{dx}=\frac{2x}{3y^2+2y-5}.\]
\end{frame}
% Diapositiva original 7
\begin{frame}{Derivación implícita}
\textbf{Ejercicio}\par\medskip Hallar la ecuación de la recta tangente a la curva $x^3+y^3=6xy$ en el punto $(3,3)$.
\end{frame}
% Diapositiva original 8
\begin{frame}{Derivadas de orden superior}
Llamamos derivada de segundo orden de una función a la derivada de la función derivada.
\[f(x)\ \longrightarrow\ f'(x)=\frac{df}{dx}\ \longrightarrow\ f''(x)=\frac{d^2f}{dx^2}.\]
La derivada tercera es la derivada de la función derivada segunda.
\[f\ \longrightarrow\ \frac{df}{dx}\ \longrightarrow\ \frac{d^2f}{dx^2}\ \longrightarrow\ \frac{d^3f}{dx^3}.\]
Y así sucesivamente se da origen a las derivadas de orden superior.
\end{frame}
% Diapositiva original 9
\begin{frame}{Derivadas de orden superior}
\small\textbf{Ejemplo}\par Obtener todas las derivadas de orden superior de $f(x)=x^5$.
\[\begin{aligned}
f(x)&=x^5,\\
f'(x)&=\frac{d}{dx}(x^5)=5x^4,\\
f''(x)&=\frac{d}{dx}(5x^4)=20x^3,\\
f^{(3)}(x)&=\frac{d}{dx}(20x^3)=60x^2,\\
f^{(4)}(x)&=\frac{d}{dx}(60x^2)=120x,\\
f^{(5)}(x)&=\frac{d}{dx}(120x)=120,\\
f^{(6)}(x)&=\frac{d}{dx}(120)=0.
\end{aligned}\]
\end{frame}
% Diapositiva original 10
\begin{frame}{Ejercicio}
Si $f(x)=\dfrac1x$, determinar $f^{(n)}(x)$.
\end{frame}
% Diapositiva original 11
\begin{frame}{Valores extremos en un intervalo}
\textbf{Máximos y mínimos relativos}\par\medskip
Sea $f$ una función definida sobre un intervalo abierto $(a,b)$ que contiene a $c$.
\begin{itemize}
\item Si $f(c)$ es un máximo de $f$ en $(a,b)$, entonces $f(c)$ es un máximo relativo.
\item Si $f(c)$ es un mínimo de $f$ en $(a,b)$, entonces $f(c)$ es un mínimo relativo.
\end{itemize}
\end{frame}
% Diapositiva original 12
\begin{frame}{Valores extremos en un intervalo}
\textbf{Puntos críticos}\par\medskip
Sea $f$ definida en $c$. Si $f'(c)=0$ o si $f$ no es derivable en $c$, entonces $c$ es un punto crítico de $f$.
\medskip

Los extremos relativos interiores ocurren sólo en los puntos críticos.
\end{frame}
% Diapositiva original 13
\begin{frame}{Funciones crecientes y decrecientes}
\textbf{Teorema}\par Sea $f$ una función continua en el intervalo cerrado $[a,b]$ y derivable en el intervalo abierto $(a,b)$.\medskip

Si $f'(x)>0$ para todo $x\in(a,b)$, $f$ es creciente en $[a,b]$.
\figura{creciente.jpg}{.90}{.48}
\end{frame}
% Diapositiva original 14
\begin{frame}{Funciones crecientes y decrecientes}
\textbf{Teorema}\par Sea $f$ una función continua en el intervalo cerrado $[a,b]$ y derivable en el intervalo abierto $(a,b)$.\medskip

Si $f'(x)<0$ para todo $x\in(a,b)$, $f$ es decreciente en $[a,b]$.
\figura{decreciente.jpg}{.90}{.48}
\end{frame}
% Diapositiva original 15
\begin{frame}{Funciones crecientes y decrecientes}
\textbf{Teorema}\par Sea $f$ una función continua en el intervalo cerrado $[a,b]$ y derivable en el intervalo abierto $(a,b)$.\medskip

Si $f'(x)=0$ para todo $x\in(a,b)$, $f$ es constante en $[a,b]$.
\figura{constante.png}{.90}{.48}
\end{frame}
% Diapositiva original 16
\begin{frame}{Funciones crecientes y decrecientes}
Si $f$ presenta un máximo o un mínimo relativo en $a$, y existe $f'(a)$, entonces $f'(a)=0$.
\medskip

En los máximos o mínimos relativos, la tangente, si existe, es horizontal y, por tanto, su derivada es cero.
\medskip

Si $f'(x)$ no fuera cero, sería positiva o negativa, por lo que $f$ sería creciente o decreciente en un entorno del punto.
\figura{extremos.jpg}{.62}{.40}
\end{frame}
% Diapositiva original 17
\begin{frame}{Criterio de la primera derivada}
\small\textbf{Teorema}\par Sea $c\in(a,b)$ un punto crítico de una función $f$ continua en $(a,b)$ y derivable en los puntos $x\ne c$.
\begin{enumerate}
\item Si $f'(x)$ cambia de negativa a positiva en $c$, $f$ tiene un mínimo relativo en $c$.
\item Si $f'(x)$ cambia de positiva a negativa en $c$, $f$ tiene un máximo relativo en $c$.
\item Si $f'(x)$ es positiva en ambos lados de $c$, o negativa en ambos lados, $f(c)$ no es ni mínimo ni máximo relativo.
\end{enumerate}
\end{frame}
% Diapositiva original 18
\begin{frame}{Funciones crecientes y decrecientes}
\small\textbf{Ejemplo}\quad $f(x)=(x^2-4)^{2/3}$.
\[f'(x)=\frac23(x^2-4)^{-1/3}(2x)=\frac{4x}{3(x^2-4)^{1/3}}.\]
\begin{columns}[c]
\begin{column}{.47\textwidth}\figura{ejemplo_extremos.png}{1}{.48}\end{column}
\begin{column}{.51\textwidth}\scriptsize
\resizebox{\textwidth}{!}{\begin{tabular}{cccl}\hline
Intervalo & Valor & Signo & Conclusión\\\hline
$(-\infty,-2)$ & $x=-3$ & $f'(-3)<0$ & Decreciente\\
$(-2,0)$ & $x=-1$ & $f'(-1)>0$ & Creciente\\
$(0,2)$ & $x=1$ & $f'(1)<0$ & Decreciente\\
$(2,\infty)$ & $x=3$ & $f'(3)>0$ & Creciente\\\hline
\end{tabular}}
\end{column}\end{columns}
\end{frame}
% Diapositiva original 19
\begin{frame}{Criterio de la segunda derivada}
\small\textbf{Teorema}\par Sea $f$ una función tal que $f'(c)=0$ y cuya segunda derivada existe en un intervalo $(a,b)$ que contiene a $c$.
\begin{enumerate}
\item Si $f''(c)>0$, $f(c)$ es un mínimo relativo.
\item Si $f''(c)<0$, $f(c)$ es un máximo relativo.
\item Si $f''(c)=0$, $f(c)$ puede o no ser un valor extremo. El criterio de la segunda derivada no permite decidir; puede aplicarse el criterio de la primera derivada.
\end{enumerate}
\end{frame}
% Diapositiva original 20
\begin{frame}{Concavidad}
Sea $f$ derivable en un intervalo $(a,b)$.
\medskip

Su gráfica es cóncava hacia arriba (convexa) si está por encima de todas sus tangentes.
\medskip

Su gráfica es cóncava hacia abajo (cóncava) si está por debajo de todas sus tangentes.
\end{frame}
% Diapositiva original 21
\begin{frame}{Concavidad}
\figura{concavidad.png}{1}{.83}
\end{frame}
% Diapositiva original 22
\begin{frame}{Concavidad}
\textbf{Criterio de concavidad}\par\medskip
\textbf{Teorema}\par Sea $f$ una función cuya segunda derivada existe en el intervalo abierto $(a,b)$.
\begin{enumerate}
\item Si $f''(x)>0$ para todo $x\in(a,b)$, $f$ es cóncava hacia arriba en $(a,b)$.
\item Si $f''(x)<0$ para todo $x\in(a,b)$, $f$ es cóncava hacia abajo en $(a,b)$.
\end{enumerate}
\end{frame}
% Diapositiva original 23
\begin{frame}{Concavidad}
\textbf{Punto de inflexión}\par\medskip
Sea $f$ una función continua en un intervalo $(a,b)$ y sea $c$ un punto de ese intervalo.
\medskip

Si la gráfica de $f$ tiene una recta tangente en $(c,f(c))$, este punto es un punto de inflexión si la concavidad de $f$ cambia de cóncava hacia arriba a cóncava hacia abajo, o viceversa, en ese punto.
\end{frame}
% Diapositiva original 24
\begin{frame}{Problemas de optimización}
La resolución del problema implica la determinación de los valores máximo y mínimo.
\end{frame}
% Diapositiva original 25
\begin{frame}{Problemas de optimización}
\small\textbf{Estrategia de resolución}
\begin{enumerate}
\item Identificar todas las cantidades dadas (datos) y las que se van a determinar (variables). Elaborar un dibujo.
\item Escribir la ecuación que se va a optimizar.
\item Reducir la ecuación a una variable independiente.
\item Determinar el dominio admisible: los intervalos de valores que tienen sentido.
\item Determinar el valor máximo o mínimo aplicando las técnicas de cálculo vistas.
\item Interpretar los resultados y rechazar los absurdos.
\end{enumerate}
\end{frame}
% Diapositiva original 26
\begin{frame}{Problemas de optimización}
\textbf{Ejemplo}\par Determinar la altura $h$ y el radio $r$ del cilindro con mayor volumen posible inscrito en una esfera de radio $a$.
\figura{cilindro.png}{.64}{.55}
\end{frame}
% Diapositiva original 27
\begin{frame}{Problemas de optimización}
\small\textbf{Ejemplo}\par Relación entre la altura $h$, el radio $r$ de la base del cilindro y el radio $a$ de la esfera, al estar inscrito el cilindro en la esfera.
\begin{columns}[c]
\begin{column}{.58\textwidth}
\[r^2=a^2-\left(\frac h2\right)^2.\]
Volumen del cilindro:
\[V=\pi r^2h.\]
\[V(h)=\pi\left(a^2-\frac{h^2}{4}\right)h.\]
\end{column}
\begin{column}{.40\textwidth}\figura{cilindro.png}{1}{.49}\end{column}
\end{columns}
\end{frame}
% Diapositiva original 28
\begin{frame}{Problemas de optimización}
\small\textbf{Ejemplo}\par La función $V(h)$ tiene sentido sólo de $0$ a $2a$:
\[V(h)=\pi\left(a^2-\frac{h^2}{4}\right)h,\quad h\in[0,2a].\]
\begin{columns}[c]
\begin{column}{.60\textwidth}
Ejemplo con $a=3$:
\[V(h)=\pi\left(9-\frac{h^2}{4}\right)h,\quad h\in[0,6].\]
\figura{volumen.png}{.90}{.37}
\end{column}
\begin{column}{.38\textwidth}\figura{cilindro.png}{1}{.45}\end{column}
\end{columns}
\end{frame}
% Diapositiva original 29
\begin{frame}{Problemas de optimización}
\small\textbf{Ejemplo}\par Localización de los puntos críticos:
\[V(h)=\pi\left(a^2-\frac{h^2}{4}\right)h,\qquad h\in[0,2a].\]
\[V'(h)=\pi\left(a^2-\frac34h^2\right).\]
\[a^2-\frac34h^2=0\quad\Longrightarrow\quad h=\frac{2a}{\sqrt3}.\]
\begin{center}\renewcommand{\arraystretch}{1.8}
\begin{tabular}{c|c|c}
$h=0$ & $h=\dfrac{2a}{\sqrt3}$ & $h=2a$\\
$V(h)=0$ & $V(h)=\dfrac{4\pi a^3}{3\sqrt3}$ & $V(h)=0$
\end{tabular}\end{center}
\end{frame}
% Diapositiva original 30
\begin{frame}{Problemas de optimización}
\small\textbf{Ejemplo}\par Sustituimos el $h$ que proporciona el $V(h)$ máximo para obtener $r$:
\[r^2=a^2-\frac{h^2}{4}=a^2-\frac{4a^2}{4\cdot3}=\frac{3a^2-a^2}{3}.\]
\begin{columns}[c]
\begin{column}{.58\textwidth}
El radio es $r=\sqrt{\dfrac23}\,a$.
\medskip

Su altura es $h=\dfrac{2a}{\sqrt3}$.
\medskip

Y su volumen $V=\dfrac{4\pi a^3}{3\sqrt3}$.
\end{column}
\begin{column}{.40\textwidth}\figura{cilindro.png}{1}{.49}\end{column}
\end{columns}
\end{frame}
% Diapositiva original 31
\begin{frame}{Problemas de optimización}
\textbf{Ejercicio}\par Hallar las dimensiones que debe tener un rectángulo inscrito en una circunferencia de $5$ cm de radio para que su área sea la mayor posible.
\figura{rectangulo.png}{.65}{.57}
\end{frame}
% Diapositiva original 32
\begin{frame}{Cálculo diferencial. Aplicaciones (I)}
\footnotesize\textbf{Ejercicios del tema}
\begin{enumerate}
\item Un fondo de inversión genera una rentabilidad $R(x)$ que depende de la cantidad de dinero invertida $x$ según
\[R(x)=-0{,}002x^2+0{,}8x-5.\]
Si tenemos 500 euros, calcular:
\begin{itemize}
\item Cuándo aumenta y cuándo disminuye la rentabilidad. Solución: aumenta en $(0,200)$ y disminuye en $(200,500)$.
\item Cuánto invertir para obtener la máxima rentabilidad y cuál será su valor. Solución: 200 euros, con rentabilidad de 75 euros.
\end{itemize}
\item Determinar los extremos locales y globales de $f(x)=3x^4-16x^3+18x^2$ en $-1\le x\le4$.
\item Hallar los máximos y mínimos locales de $g(x)=x+2\sin x$, $0\le x\le2\pi$.
\end{enumerate}
\end{frame}
% Diapositiva original 33
\begin{frame}{Cálculo diferencial. Aplicaciones (I)}
\footnotesize
\textbf{4)} Determinar la derivada segunda de $x^4+y^4=16$.
\medskip

\textbf{5)} Una empresa de fabricación de puertas de madera utiliza un tablón rectangular para la hoja y tres listones de 10 cm de ancho para el marco (dos laterales y el lado superior). El precio del tablón es de 128 euros por metro cuadrado y el de los listones, 87 euros por metro lineal. Calcular:
\begin{columns}[c]
\begin{column}{.72\textwidth}
\begin{enumerate}
\item Las dimensiones de una puerta de $2\,\mathrm{m}^2$ de superficie de hoja para que el coste sea mínimo y su precio. Solución: $2\times1$ m; 621,4 euros.
\item Si la puerta tiene 2,5 m de ancho y 0,8 m de alto, su precio. Solución: 630,1 euros.
\end{enumerate}
\end{column}
\begin{column}{.26\textwidth}\figura{puerta.png}{1}{.35}\end{column}
\end{columns}
\end{frame}
\end{document}
